\chapter{Conclusion}
\label{chap:conclusion}

% Summary of what the project achieved relative to its stated objectives
\section{Achievements}
\label{sec:achievements}

The main project achievements are summarised below:

\begin{itemize}
  \item Traffic-demand data was converted into vehicle flows suitable for the SUMO simulation, including PCU-based vehicle classes and OD allocation.
  \item The Thapathali intersection was configured in SUMO with road layout, traffic flow, and signal-phase representation.
  \item A fixed 60-second signal controller was implemented as the baseline for comparison.
  \item A DQN-based controller was designed and trained for adaptive signal phase selection.
  \item On the evaluated 6 AM to 10 PM period, average waiting time was reduced by 22.90\% compared with the fixed baseline.
  \item Average queue length was reduced by 30.06\%, and throughput per waiting second improved by 32.65\%.
\end{itemize}

% Constraints, shortcomings, and areas where further work is warranted
\section{Limitations}
\label{sec:limitations}

The results should be interpreted with the following limitations in mind:

\begin{itemize}
  \item The controller was tested only in simulation; no hardware signal controller or field trial was implemented.
  \item Pedestrians and public transport were excluded to keep the control problem tractable and focused on vehicle-delay reduction at Thapathali.
  \item Unusual traffic conditions such as accidents, road closures, festivals, or sudden diversions were not represented in the training data.
  \item Random traffic signal timers were placed near outgoing lanes to approximate downstream backspill, but they were not calibrated using actual neighbouring-intersection data.
  \item The study considers one intersection only, so it does not capture network-wide coordination effects.
\end{itemize}

% Directions for extending this work in scope, methodology, or deployment
\section{Future Work}
\label{sec:future_work}

Future work can extend the project in the following directions:

\begin{itemize}
  \item \textbf{Multi-intersection coordination:} Extend the controller to nearby connected intersections so that one junction's discharge pattern does not create unmanaged pressure downstream.
  \item \textbf{Faster training experiments:} Compare SUMO with lighter simulators such as CityFlow to test more episodes and demand variations, while noting the trade-off between speed and physical realism.
  \item \textbf{Richer traffic modes:} Add pedestrians, buses, public transport stops, and non-motorised traffic so that the controller reflects more of the real Thapathali operating environment.
  \item \textbf{Robustness testing:} Evaluate the trained policy under abnormal demand patterns such as sudden surges, lane blockage, and event-day traffic.
\end{itemize}

% Closing statement on the project's overall contribution and significance
\section{Final Remarks}
\label{sec:final_remarks}

This project set out to investigate whether a reinforcement
learning-based approach could meaningfully improve traffic signal
control at the Thapathali intersection in Kathmandu, one of the
city\textquotesingle s most persistently congested arterial junctions.
The results suggest that it can, and the case for adaptive signal
control is stronger than the numbers alone convey.

The DQN agent demonstrated consistent improvement over fixed-timer logic
across the full operational window. Where the gains are most evident is
during low-to-moderate demand periods, hours where fixed timing is
structurally wasteful, holding vehicles at red signals long after
opposing flows have cleared. In these windows, the
agent\textquotesingle s ability to read and react to real-time queue
states yielded noticeably better outcomes, highlighting just how rigid
and inefficient schedule-based control becomes when traffic is light and
variable.

The more telling story, however, lies in the peak hours. Improvements
during high-demand periods, the morning and evening surges that define
Kathmandu\textquotesingle s daily congestion profile, were comparatively
modest. This is not a failure of the algorithm; it is a reflection of a
harder problem. At genuine saturation, when every approach is
simultaneously overloaded, no signal timing strategy alone can conjure
capacity that does not exist. The Thapathali intersection, like many in
the Kathmandu valley, operates well beyond its designed throughput
during peak windows. The physical geometry, lane discipline, and
competing turning movements create a level of congestion that signal
optimisation can soften but cannot resolve. The DQN agent improves the
allocation of insufficient green time; it does not manufacture more road
space.

This distinction matters for how the results should be interpreted. The
system performs well precisely where fixed timers fail most obviously,
and it reaches its natural ceiling precisely where the underlying
infrastructure is the binding constraint. That ceiling is not an
algorithmic limitation but an infrastructural one, and it points toward
what complementary interventions would be needed alongside smarter
signal control to meaningfully address peak-hour gridlock.

As a proof of concept, this work establishes a credible and reproducible
foundation. The simulation pipeline, agent design, and evaluation
framework are all in place. What lies ahead, including real-world
hardware integration, inclusion of pedestrians and public transport, and
coordination across multiple connected intersections, are well-defined
extensions rather than open questions, each building directly on what
has been demonstrated here.

Urban traffic management in cities like Kathmandu has long depended on
static, schedule-driven systems that cannot adapt to the variability and
intensity of real-world demand. This project offers substantive early
evidence that adaptive, learning-based control can meaningfully close
that gap, not as a complete solution, but as an important and scalable
component of one.
